\documentclass[12pt,letterpaper]{article}
\usepackage{graphicx,textcomp}
\usepackage{natbib}
\usepackage{setspace}
\usepackage[margin=1in]{geometry}
\usepackage{color}
\usepackage[reqno]{amsmath}
\usepackage{amsthm}
\usepackage{fancyvrb}
\usepackage{amssymb,enumerate}
\usepackage[all]{xy}
\usepackage{endnotes}
\usepackage{lscape}
\newtheorem{com}{Comment}
\usepackage{float}
\usepackage{hyperref}
\newtheorem{lem} {Lemma}
\newtheorem{prop}{Proposition}
\newtheorem{thm}{Theorem}
\newtheorem{defn}{Definition}
\newtheorem{cor}{Corollary}
\newtheorem{obs}{Observation}
\usepackage[compact]{titlesec}
\usepackage{dcolumn}
\usepackage{tikz}
\usetikzlibrary{arrows}
\usepackage{multirow}
\usepackage{xcolor}
\newcolumntype{.}{D{.}{.}{-1}}
\newcolumntype{d}[1]{D{.}{.}{#1}}
\definecolor{light-gray}{gray}{0.65}
\usepackage{url}
\usepackage{listings}
\usepackage{color}

\definecolor{codegreen}{rgb}{0,0.6,0}
\definecolor{codegray}{rgb}{0.5,0.5,0.5}
\definecolor{codepurple}{rgb}{0.58,0,0.82}
\definecolor{backcolour}{rgb}{0.95,0.95,0.92}

\lstdefinestyle{mystyle}{
	backgroundcolor=\color{backcolour},   
	commentstyle=\color{codegreen},
	keywordstyle=\color{magenta},
	numberstyle=\tiny\color{codegray},
	stringstyle=\color{codepurple},
	basicstyle=\footnotesize,
	breakatwhitespace=false,         
	breaklines=true,                 
	captionpos=b,                    
	keepspaces=true,                 
	numbers=left,                    
	numbersep=5pt,                  
	showspaces=false,                
	showstringspaces=false,
	showtabs=false,                  
	tabsize=2
}
\lstset{style=mystyle}
\newcommand{\Sref}[1]{Section~\ref{#1}}
\newtheorem{hyp}{Hypothesis}


\title{Problem Set 4}
\date{Due: November 27, 2025}
\author{Applied Stats/Quant Methods 1}


\begin{document}
	\maketitle
	\section*{Instructions}
	\begin{itemize}
		\item Please show your work! You may lose points by simply writing in the answer. If the problem requires you to execute commands in \texttt{R}, please include the code you used to get your answers. Please also include the \texttt{.R} file that contains your code. If you are not sure if work needs to be shown for a particular problem, please ask.
		\item Your homework should be submitted electronically on GitHub.
		\item This problem set is due before 23:59 on Thursday November 27, 2025. No late assignments will be accepted.
	\end{itemize}



	\vspace{.5cm}
\section*{Question 1: Economics}
\vspace{.25cm}
\noindent 	
In this question, use the \texttt{prestige} dataset in the \texttt{car} library. First, run the following commands:

\begin{verbatim}
install.packages(car)
library(car)
data(Prestige)
help(Prestige)
\end{verbatim} 


\noindent We would like to study whether individuals with higher levels of income have more prestigious jobs. Moreover, we would like to study whether professionals have more prestigious jobs than blue and white collar workers.

\newpage
\begin{enumerate}
	
	\item [(a)]
	Create a new variable \texttt{professional} by recoding the variable \texttt{type} so that professionals are coded as $1$, and blue and white collar workers are coded as $0$ (Hint: \texttt{ifelse}).
	
	\lstinputlisting[language=R, firstline=39, lastline=41]{PS04_DP.R}
	\vspace{0.5cm}
	
	
	\item [(b)]
	Run a linear model with \texttt{prestige} as an outcome and \texttt{income}, \texttt{professional}, and the interaction of the two as predictors (Note: this is a continuous $\times$ dummy interaction.)
	\lstinputlisting[language=R, firstline=43, lastline=46]{PS04_DP.R}
\begin{table}[!htbp] \centering 
	\caption{} 
	\label{} 
	\begin{tabular}{@{\extracolsep{5pt}}lc} 
		\\[-1.8ex]\hline 
		\hline \\[-1.8ex] 
		& \multicolumn{1}{c}{\textit{Dependent variable:}} \\ 
		\cline{2-2} 
		\\[-1.8ex] & prestige \\ 
		\hline \\[-1.8ex] 
		income & 0.003$^{***}$ \\ 
		& (0.0005) \\ 
		& \\ 
		professional & 37.781$^{***}$ \\ 
		& (4.248) \\ 
		& \\ 
		interaction term & $-$0.002$^{***}$ \\ 
		& (0.001) \\ 
		& \\ 
		Constant & 21.142$^{***}$ \\ 
		& (2.804) \\ 
		& \\ 
		\hline \\[-1.8ex] 
		\textit{Note:}  & \multicolumn{1}{r}{$^{*}$p$<$0.1; $^{**}$p$<$0.05; $^{***}$p$<$0.01} \\ 
	\end{tabular} 
\end{table} 
	
	\vspace{0.5cm}
	\item [(c)]
	Write the prediction equation based on the result.
		\begin{center} 
		$\widehat {\text{Prestige}} = \hat{\alpha} + \hat{\beta_1} \times \text{Income} + \hat{\beta_2} \times \text{Professional} + \hat{\beta_3} \times \text{Income} \times \text{Professional} $ \\[0.2cm]
		$\widehat {\text{Prestige}} = 21.142 + 0.003 \times \text{Income} + 37.781 \times \text{Professional} - 0.002\times \text{Income} \times \text{Professional} $
	\end{center}
	
\newpage
	\item [(d)]
	Interpret the coefficient for \texttt{income}. \\[0.2cm]
	For non-professionals, for every one-unit increase in \texttt{income}, there is a 0.003 unit increase in \texttt{prestige}, on average. 
	\vspace{0.5cm}	
	\item [(e)]
	Interpret the coefficient for \texttt{professional}. \\[0.2cm]
	On average, when \texttt{income} is held at zero, professional occupations are associated with a 37.781-point increase in prestige compared to non-professional occupations.
	\vspace{0.5cm}	
	\item [(f)]
	What is the effect of a \$1,000 increase in income on prestige score for professional occupations? In other words, we are interested in the marginal effect of income when the variable \texttt{professional} takes the value of $1$. Calculate the change in $\hat{y}$ associated with a \$1,000 increase in income based on your answer for (c).
	\lstinputlisting[language=R, firstline=48, lastline=52]{PS04_DP.R}
	\begin{verbatim}
		0.8452 
	\end{verbatim}
	To find the effect of a \$1,000 increase in income for professional occupations, the predicted prestige score at \texttt{income = 0} was subtracted from the predicted prestige score at \texttt{income = 1000}, holding \texttt{professional = 1} in both cases. Plugging the values into the prediction equation, the difference between the two is 0.8452. Thus, the change in \texttt{prestige} associated with a \$1,000 increase in income is 0.8452 units. 
	\vspace{0.5cm}
	\item [(g)]
	What is the effect of changing one's occupations from non-professional to professional when her income is \$6,000? We are interested in the marginal effect of professional jobs when the variable \texttt{income} takes the value of $6,000$. Calculate the change in $\hat{y}$ based on your answer for (c).
	\lstinputlisting[language=R, firstline=54, lastline=58]{PS04_DP.R}
	\begin{verbatim}
		23.82703
	\end{verbatim}
	To find the effect of switching from a non-professional to a professional occupation when income is \$6,000, the predicted prestige score at \texttt{professional = 0} was subtracted from the score found at \texttt{professional = 1}, holding \texttt{income = 6000} in both cases. Plugging the values into the prediction equation, the difference between the two is 23.82703. Therefore, when income is \$6,000, changing occupation from a non-professional one to a professional one increases the predicted prestige score by approximately 23.83 units. 
\newpage
\end{enumerate}
\section*{Question 2: Political Science}
\vspace{.25cm}
\noindent 	Researchers are interested in learning the effect of all of those yard signs on voting preferences.\footnote{Donald P. Green, Jonathan	S. Krasno, Alexander Coppock, Benjamin D. Farrer,	Brandon Lenoir, Joshua N. Zingher. 2016. ``The effects of lawn signs on vote outcomes: Results from four randomized field experiments.'' Electoral Studies 41: 143-150. } Working with a campaign in Fairfax County, Virginia, 131 precincts were randomly divided into a treatment and control group. In 30 precincts, signs were posted around the precinct that read, ``For Sale: Terry McAuliffe. Don't Sellout Virgina on November 5.'' \\

Below is the result of a regression with two variables and a constant.  The dependent variable is the proportion of the vote that went to McAuliff's opponent Ken Cuccinelli. The first variable indicates whether a precinct was randomly assigned to have the sign against McAuliffe posted. The second variable indicates
a precinct that was adjacent to a precinct in the treatment group (since people in those precincts might be exposed to the signs).  \\

\vspace{.5cm}
\begin{table}[!htbp]
	\centering 
	\textbf{Impact of lawn signs on vote share}\\
	\begin{tabular}{@{\extracolsep{5pt}}lccc} 
		\\[-1.8ex] 
		\hline \\[-1.8ex]
		Precinct assigned lawn signs  (n=30)  & 0.042\\
		& (0.016) \\
		Precinct adjacent to lawn signs (n=76) & 0.042 \\
		&  (0.013) \\
		Constant  & 0.302\\
		& (0.011)
		\\
		\hline \\
	\end{tabular}\\
	\footnotesize{\textit{Notes:} $R^2$=0.094, N=131}
\end{table}

\vspace{.5cm}
\begin{enumerate}
	\item [(a)] Use the results from a linear regression to determine whether having these yard signs in a precinct affects vote share (e.g., conduct a hypothesis test with $\alpha = .05$). \\ [0.5 cm]
		\textbf{Hypotheses} \\[0.3cm]
		$ H_{0}: \beta_i= 0; H_{A}: \beta_i\neq 0$ \\ [0.3cm]
		\textbf{Test Statistic}
		\lstinputlisting[language=R, firstline=65, lastline=66]{PS04_DP.R}
		\begin{verbatim}
			[1] 2.625
		\end{verbatim}
	\newpage		
		\textbf{Calculate the p-value}
	\lstinputlisting[language=R, firstline=68, lastline=73]{PS04_DP.R}
	\begin{verbatim}
		[1] 0.00972002
	\end{verbatim}
		\textbf{Conclusion} \\[0.3cm]
		The p-value is smaller than 0.05, so we can reject the null hypothesis that the coefficient for the precinct assigned lawn signs is 0. Thus, being in a precinct with yard signs has a statistically significant effect on vote share. 
	\item [(b)]  Use the results to determine whether being
	next to precincts with these yard signs affects vote
	share (e.g., conduct a hypothesis test with $\alpha = .05$). \\[0.5cm]
	\textbf{Hypotheses} \\ [0.3cm]
	$ H_{0}: \beta_j= 0; H_{A}: \beta_j\neq 0$ \\ [0.3cm]
	\textbf{Test Statistic}
	\lstinputlisting[language=R, firstline=79, lastline=80]{PS04_DP.R}
	\begin{verbatim}
		[1] 3.230769
	\end{verbatim}
	\textbf{Calculate the p-value}
	\lstinputlisting[language=R, firstline=82, lastline=84]{PS04_DP.R}
	\begin{verbatim}
		[1] 0.00156946
	\end{verbatim}
		\textbf{Conclusion} \\[0.3cm]
		The p-value is smaller than 0.05, so we can reject the null hypothesis that the coefficient for the precincts adjacent to lawn signs is 0. Thus, proximity to precincts with yard signs has a statistically significant effect on vote share. 
	
	\item [(c)] Interpret the coefficient for the constant term substantively. \\[0.5cm]
	The constant (intercept) represents the predicted vote share for Ken Cuccinelli in precincts where no lawn signs were posted and that were not adjacent to these precincts. Thus, the baseline predicted vote share for Cuccinelli in this type of untreated precincts is 30.2\%.
	\item [(d)] Evaluate the model fit for this regression.  What does this	tell us about the importance of yard signs versus other factors that are not modeled? \\[0.5cm]
	The model has an $R^2$ value of 0.094, so it only explains 9.4\% of the variation in the outcome variable, in the vote share for Cuccinelli. While it was shown above that the presence of lawn signs in a precinct or the proximity to them has a statistically significant effect on the vote share, most of the variation in outcome is explained by other factors, not present in this model. Lawn signs have some effect on vote share, but they don't explain much of its overall variation. 
	
\end{enumerate}  


\end{document}
